\documentclass[12pt,a4paper]{article}

\usepackage[utf8]{inputenc}
\usepackage[T1]{fontenc}
\usepackage[brazil]{babel}
\usepackage[top=3cm,left=3cm,bottom=2cm,right=2cm]{geometry}
\usepackage{mathptmx}
\usepackage{inconsolata}
\usepackage{setspace}
\usepackage{indentfirst}
\usepackage{titlesec}
\usepackage{booktabs}
\usepackage{array}
\usepackage{listings}
\usepackage{xcolor}
\usepackage{tikz}

\onehalfspacing
\setlength{\parindent}{1.25cm}

\titleformat*{\section}{\normalfont\bfseries\normalsize}
\titleformat*{\subsection}{\normalfont\bfseries\normalsize}

\lstset{
    basicstyle=\ttfamily\footnotesize,
    breaklines=true,
    columns=fullflexible,
    keepspaces=true,
    frame=single,
    rulecolor=\color{black!25},
    backgroundcolor=\color{black!4},
    xleftmargin=0.2cm,
    xrightmargin=0.2cm,
    framexleftmargin=4pt,
    framexrightmargin=4pt,
    framextopmargin=4pt,
    framexbottommargin=4pt,
    literate=
        {á}{{\'a}}1 {é}{{\'e}}1 {í}{{\'i}}1 {ó}{{\'o}}1 {ú}{{\'u}}1
        {Á}{{\'A}}1 {É}{{\'E}}1 {Í}{{\'I}}1 {Ó}{{\'O}}1 {Ú}{{\'U}}1
        {ã}{{\~a}}1 {õ}{{\~o}}1 {Ã}{{\~A}}1 {Õ}{{\~O}}1
        {â}{{\^a}}1 {ê}{{\^e}}1 {ô}{{\^o}}1
        {à}{{\`a}}1 {ç}{{\c{c}}}1 {Ç}{{\c{C}}}1
}

\begin{document}

\section{Questão 06: Controle fuzzy da ventilação de uma sala}

O sistema recebe a temperatura e a umidade relativa de uma sala e calcula o comando percentual de um ventilador. As regras favorecem ventilação mais alta em condições quentes ou úmidas e produzem transições graduais. O modelo foi implementado em Python com a inferência de Mamdani [1].

\subsection{Variáveis e funções de pertinência}

As entradas são temperatura $T$ em $[16, 40]\;^{\circ}$C e umidade relativa $U$ em $[0, 100]\%$. A saída $V$ tem universo $[0, 100]\%$ e representa o comando do ventilador. Os intervalos e os pontos da tabela são escolhas didáticas e podem ser ajustados para outra sala.

\begin{table}[htbp]
    \centering
    \caption{Pontos $(a, b, c, d)$ das funções de pertinência}
    \begin{tabular}{>{\raggedright\arraybackslash}p{4.2cm} >{\raggedright\arraybackslash}p{2.5cm} >{\raggedright\arraybackslash}p{4.2cm}}
        \toprule
        \textbf{Variável} & \textbf{Termo} & \textbf{Pontos (a, b, c, d)} \\
        \midrule
        Temperatura ($^{\circ}$C) & Fria & (16, 16, 20, 25) \\
        \addlinespace
        Temperatura ($^{\circ}$C) & Amena & (21, 26, 26, 31) \\
        \addlinespace
        Temperatura ($^{\circ}$C) & Quente & (27, 32, 40, 40) \\
        \addlinespace
        Umidade (\%) & Baixa & (0, 0, 30, 50) \\
        \addlinespace
        Umidade (\%) & Média & (35, 55, 55, 75) \\
        \addlinespace
        Umidade (\%) & Alta & (60, 80, 100, 100) \\
        \addlinespace
        Ventilação (\%) & Baixa & (0, 20, 20, 40) \\
        \addlinespace
        Ventilação (\%) & Média & (25, 50, 50, 75) \\
        \addlinespace
        Ventilação (\%) & Alta & (60, 80, 80, 100) \\
        \bottomrule
    \end{tabular}
\end{table}

Cada função vale zero fora de $[a, d]$, cresce linearmente de 0 a 1 entre $a$ e $b$, vale 1 entre $b$ e $c$ e decresce linearmente de 1 a 0 entre $c$ e $d$. Quando $b = c$, a função é triangular. Quando $a = b$ ou $c = d$, a rampa correspondente desaparece e a pertinência vale 1 na extremidade. A sobreposição permite que uma leitura pertença parcialmente a mais de um termo.

\subsection{Base de regras}

Cada célula representa uma regra. Por exemplo, SE a temperatura é amena E a umidade é alta, ENTÃO a ventilação é alta. As nove combinações cobrem todas as entradas válidas.

\begin{table}[htbp]
    \centering
    \begin{tabular}{>{\raggedright\arraybackslash}p{4.8cm} >{\raggedright\arraybackslash}p{2.2cm} >{\raggedright\arraybackslash}p{2.2cm} >{\raggedright\arraybackslash}p{2.2cm}}
        \toprule
        \textbf{Temperatura / Umidade} & \textbf{Baixa} & \textbf{Média} & \textbf{Alta} \\
        \midrule
        Fria & Baixa & Baixa & Média \\
        \addlinespace
        Amena & Baixa & Média & Alta \\
        \addlinespace
        Quente & Média & Alta & Alta \\
        \bottomrule
    \end{tabular}
\end{table}

\subsection{Inferência de Mamdani e defuzzificação}

Primeiro, o sistema calcula os graus de pertinência das duas entradas. Em cada regra, o operador E usa o mínimo dos dois graus. Esse valor limita a altura do conjunto fuzzy consequente. O máximo ponto a ponto combina os consequentes cortados. Se várias regras ativam o mesmo termo de saída, basta usar o maior corte desse termo.

O centroide divide o momento da curva agregada em relação à origem pela sua área. O programa aproxima as duas integrais pela regra dos trapézios, com passo de 0,1 ponto percentual entre 0 e 100. O resultado é um comando numérico para o ventilador.

\subsection{Exemplo com 30 $^{\circ}$C e 65\% de umidade}

A temperatura tem pertinência 0 em fria, $(31 - 30)/(31 - 26) = 0,20$ em amena e $(30 - 27)/(32 - 27) = 0,60$ em quente. A umidade tem pertinência 0 em baixa, $(75 - 65)/(75 - 55) = 0,50$ em média e $(65 - 60)/(80 - 60) = 0,25$ em alta.

\begin{table}[htbp]
    \centering
    \begin{tabular}{>{\raggedright\arraybackslash}p{2.8cm} >{\raggedright\arraybackslash}p{2.5cm} >{\raggedright\arraybackslash}p{2.8cm} >{\raggedright\arraybackslash}p{3.8cm}}
        \toprule
        \textbf{Temperatura} & \textbf{Umidade} & \textbf{Ventilação} & \textbf{Ativação (mínimo)} \\
        \midrule
        Amena & Média & Média & 0,20 \\
        \addlinespace
        Amena & Alta & Alta & 0,20 \\
        \addlinespace
        Quente & Média & Alta & 0,50 \\
        \addlinespace
        Quente & Alta & Alta & 0,25 \\
        \bottomrule
    \end{tabular}
\end{table}

A ventilação média fica cortada em 0,20. Para a ventilação alta, o maior corte é 0,50. A baixa fica em zero. O gráfico mostra a união dos conjuntos cortados e a posição do centroide.

\begin{figure}[htbp]
    \centering
    \begin{tikzpicture}[x=0.105cm, y=3.0cm]
        \fill[gray!20] (25,0) -- (30,0.2) -- (64,0.2) -- (70,0.5) -- (90,0.5) -- (100,0) -- cycle;
        \draw[thick] (0,0) -- (25,0) -- (30,0.2) -- (64,0.2) -- (70,0.5) -- (90,0.5) -- (100,0);
        \draw (0,0) -- (100,0);
        \draw (0,0) -- (0,1.0);
        \foreach \x in {0,20,40,60,80,100} {
            \draw (\x,0) -- (\x,-0.04) node[below] {\small \x};
        }
        \foreach \y/\ylab in {0.2/{0,2}, 0.5/{0,5}, 1.0/{1,0}} {
            \draw (0,\y) -- (-1.2,\y) node[left] {\small \ylab};
        }
        \node[above right] at (0,1.02) {\small Grau de pertinência};
        \node at (50,-0.34) {\small Comando de ventilação (\%)};
        \draw[dashed, thick] (68.89,0) -- (68.89,0.72) node[above] {\small 68,89\%};
    \end{tikzpicture}
\end{figure}

A área aproximada é 21,9000 e o momento é 1508,7000. A razão 1508,7000 / 21,9000 resulta em um comando de \textbf{68,89\%}.

\subsection{Execução e verificação}

Execute \texttt{python ventilacao\_fuzzy.py 30 65} na pasta da questão.

O controle retornou 20\% a 16 $^{\circ}$C e 0\%, 50\% a 26 $^{\circ}$C e 55\%, e 80\% a 40 $^{\circ}$C e 100\%.

A saída pode diminuir em algumas transições. Com umidade fixa em 38\%, elevar a temperatura de 23 $^{\circ}$C para 24 $^{\circ}$C reduz o comando de 30,34\% para 28,46\%. As ativações dos consequentes variam em ritmos diferentes na sobreposição, o que desloca o centroide da saída agregada.

\section*{Referências}

\noindent [1] MAMDANI, E. H.; ASSILIAN, S. \textbf{An Experiment in Linguistic Synthesis with a Fuzzy Logic Controller}. 1975. Disponível em: https://doi.org/10.1016/S0020-7373(75)80002-2.

\end{document}
